\documentclass[12pt,a4paper]{article}

\usepackage[T2A]{fontenc}
\usepackage[utf8]{inputenc}
\usepackage[russian]{babel}
\usepackage{amsmath,amssymb,amsthm}
\usepackage{array}
\usepackage{booktabs}
\usepackage{enumitem}
\usepackage{listings}
\usepackage{xcolor}
\usepackage{geometry}

\geometry{margin=2cm}

\newtheorem{definition}{Определение}
\newtheorem{theorem}{Теорема}
\newtheorem{example}{Пример}
\newtheorem{remark}{Замечание}

\lstset{
    basicstyle=\ttfamily\small,
    frame=single,
    breaklines=true,
    columns=fullflexible,
    keywordstyle=\color{blue},
    commentstyle=\color{gray}
}

\title{Билет №40 \\ \small Отладка, тестирование, верификация и оценивание сложности программ. Генерация тестов. Системы генерации тестов. (ИТМО) \\ Методы тестирования программного обеспечения и области их применения. (СпБПУ)}
\author{Сериков Владислав Александрович}
\date{13 мая 2026}

\begin{document}

\maketitle
\tableofcontents
\newpage

\section{Базовые понятия качества программного обеспечения}

\begin{definition}
\textbf{Программная ошибка} --- это несоответствие программы требованиям, спецификации, ожиданиям пользователя или внутренним инвариантам системы.
\end{definition}

Обычно различают следующие понятия:

\begin{itemize}
    \item \textbf{mistake} --- ошибка человека, например неправильное понимание требования;
    \item \textbf{defect, fault, bug} --- дефект в артефакте: коде, модели, документации, тесте;
    \item \textbf{error} --- неправильное внутреннее состояние программы;
    \item \textbf{failure} --- наблюдаемое пользователем неправильное поведение программы.
\end{itemize}

\begin{definition}
\textbf{Тестирование программного обеспечения} --- это процесс выполнения программы или анализа её артефактов с целью обнаружения дефектов и получения информации о качестве продукта.
\end{definition}

Тестирование не доказывает отсутствие ошибок в общем случае. Оно показывает наличие дефектов, а также повышает уверенность в корректности программы относительно выбранной модели требований и выбранного набора тестов.

\begin{definition}
\textbf{Отладка} --- это процесс локализации, анализа и устранения причины обнаруженного дефекта.
\end{definition}

\begin{definition}
\textbf{Верификация} --- проверка того, что программный продукт построен правильно относительно формальной или неформальной спецификации. Кратко: ``строим ли мы продукт правильно?''
\end{definition}

\begin{definition}
\textbf{Валидация} --- проверка того, что программный продукт удовлетворяет реальным потребностям пользователя. Кратко: ``правильный ли продукт мы строим?''
\end{definition}

\section{Отладка программ}

\subsection{Цели и этапы отладки}

Отладка начинается после обнаружения сбоя, например после падения теста, аварийного завершения программы, нарушения инварианта или жалобы пользователя.

Типовой процесс отладки:

\begin{enumerate}
    \item \textbf{Воспроизведение дефекта.}
    Нужно получить минимальный сценарий, при котором ошибка проявляется стабильно.

    \item \textbf{Сбор информации.}
    Используются логи, трассировка, дампы памяти, профилировщики, отладчики, системные метрики.

    \item \textbf{Локализация причины.}
    Определяется участок кода, состояние или входные данные, приводящие к ошибке.

    \item \textbf{Построение гипотезы.}
    Формулируется предположение о причине дефекта.

    \item \textbf{Проверка гипотезы.}
    Гипотеза проверяется экспериментом: дополнительным логированием, тестом, анализом переменных.

    \item \textbf{Исправление.}
    Вносится минимальное изменение, устраняющее причину, а не только симптом.

    \item \textbf{Регрессионное тестирование.}
    Проверяется, что дефект исправлен и не появились новые ошибки.

    \item \textbf{Документирование.}
    Фиксируются причина, исправление, тесты и возможные меры профилактики.
\end{enumerate}

\subsection{Основные методы отладки}

\subsubsection{Отладка с помощью вывода}

Простейший метод --- добавить в программу диагностический вывод.

\begin{lstlisting}[language=C]
printf("x = %d, y = %d\n", x, y);
\end{lstlisting}

Преимущества:
\begin{itemize}
    \item простота;
    \item не требуется специальный инструмент;
    \item удобно для серверных и распределённых приложений.
\end{itemize}

Недостатки:
\begin{itemize}
    \item избыточный вывод затрудняет анализ;
    \item логирование может менять временное поведение программы;
    \item после исправления часть диагностического кода нужно удалить или скрыть за уровнем логирования.
\end{itemize}

\subsubsection{Пошаговая отладка}

Используется интерактивный отладчик: \texttt{gdb}, \texttt{lldb}, отладчики IDE.

Основные операции:

\begin{itemize}
    \item установка точки останова;
    \item пошаговое выполнение;
    \item просмотр значений переменных;
    \item анализ стека вызовов;
    \item просмотр памяти;
    \item условные точки останова;
    \item watchpoint --- остановка при изменении переменной или области памяти.
\end{itemize}

\subsubsection{Дельта-отладка}

\begin{definition}
\textbf{Дельта-отладка} --- это метод автоматического уменьшения входных данных или изменений программы до минимального набора, всё ещё вызывающего сбой.
\end{definition}

\textbf{Идея:} если большой входной файл вызывает ошибку, то можно попытаться удалить из него части и проверить, сохраняется ли сбой.

\paragraph{Алгоритм упрощённой дельта-отладки входа}

Пусть есть вход $I$, на котором программа падает.

\begin{enumerate}
    \item Разбить $I$ на $n$ частей.
    \item Для каждой части попробовать удалить её.
    \item Если после удаления сбой сохраняется, заменить $I$ уменьшенным входом.
    \item Если ни одну часть удалить нельзя, увеличить детализацию разбиения.
    \item Повторять, пока дальнейшее уменьшение невозможно.
\end{enumerate}

\begin{example}
Пусть парсер падает на строке:
\[
\texttt{begin a = 1; b = 2; c = (3 + ); end}
\]
Дельта-отладка может последовательно уменьшить её до:
\[
\texttt{c=(3+);}
\]
Это минимальный фрагмент, демонстрирующий ошибку обработки незавершённого выражения.
\end{example}

\subsubsection{Отладка по инвариантам}

В программу добавляют утверждения, которые должны быть истинны в определённых точках.

\begin{lstlisting}[language=C]
assert(size >= 0);
assert(index < size);
\end{lstlisting}

Инварианты позволяют обнаружить ошибку ближе к моменту её возникновения, а не в точке позднего сбоя.

\subsection{Типичные причины дефектов}

\begin{itemize}
    \item неверная обработка граничных случаев;
    \item ошибки индексации массивов;
    \item переполнение целых чисел;
    \item деление на ноль;
    \item ошибки работы с памятью;
    \item состояние гонки в многопоточных программах;
    \item неправильная синхронизация;
    \item нарушение протокола взаимодействия компонентов;
    \item неверное понимание требований;
    \item отсутствие обработки исключительных ситуаций.
\end{itemize}

\section{Тестирование программного обеспечения}

\subsection{Цели тестирования}

Тестирование применяется для:

\begin{itemize}
    \item обнаружения дефектов;
    \item подтверждения соответствия требованиям;
    \item оценки качества и надёжности;
    \item предотвращения регрессий;
    \item получения информации о рисках выпуска;
    \item проверки нефункциональных характеристик: производительности, безопасности, удобства использования.
\end{itemize}

\subsection{Принципы тестирования}

\begin{enumerate}
    \item Тестирование показывает наличие дефектов, но не доказывает их отсутствие.
    \item Исчерпывающее тестирование обычно невозможно.
    \item Раннее тестирование снижает стоимость исправления дефектов.
    \item Дефекты часто группируются в отдельных модулях.
    \item Одни и те же тесты со временем перестают находить новые дефекты.
    \item Тестирование зависит от контекста.
    \item Отсутствие найденных ошибок не означает полезности системы, если она не решает задачу пользователя.
\end{enumerate}

\subsection{Уровни тестирования}

\begin{center}
\begin{tabular}{p{4cm} p{9cm}}
\toprule
\textbf{Уровень} & \textbf{Назначение} \\
\midrule
Модульное тестирование & Проверка отдельных функций, классов, модулей. \\
Интеграционное тестирование & Проверка взаимодействия компонентов. \\
Системное тестирование & Проверка всей системы как единого продукта. \\
Приёмочное тестирование & Проверка соответствия ожиданиям заказчика или пользователя. \\
Регрессионное тестирование & Проверка, что изменения не сломали ранее работавшую функциональность. \\
\bottomrule
\end{tabular}
\end{center}

\subsection{Виды тестирования по доступности кода}

\subsubsection{Тестирование чёрного ящика}

\begin{definition}
\textbf{Тестирование чёрного ящика} --- тестирование без знания внутренней реализации программы, только по входам, выходам и спецификации.
\end{definition}

Применяется при:
\begin{itemize}
    \item функциональном тестировании;
    \item приёмочном тестировании;
    \item тестировании API;
    \item тестировании пользовательского интерфейса;
    \item проверке требований.
\end{itemize}

\subsubsection{Тестирование белого ящика}

\begin{definition}
\textbf{Тестирование белого ящика} --- тестирование с использованием информации о внутренней структуре программы: кода, графа потока управления, ветвлений, условий.
\end{definition}

Применяется при:
\begin{itemize}
    \item модульном тестировании;
    \item анализе покрытия кода;
    \item проверке сложной логики;
    \item тестировании критичных участков;
    \item анализе безопасности.
\end{itemize}

\subsubsection{Тестирование серого ящика}

\begin{definition}
\textbf{Тестирование серого ящика} --- тестирование, при котором тестировщик частично знает внутреннюю структуру системы.
\end{definition}

Применяется при интеграционном тестировании, тестировании веб-приложений, API и микросервисов.

\section{Методы функционального тестирования}

\subsection{Разбиение на классы эквивалентности}

\begin{definition}
\textbf{Класс эквивалентности} --- множество входных данных, которые с точки зрения спецификации должны обрабатываться одинаково.
\end{definition}

\paragraph{Алгоритм построения тестов}

\begin{enumerate}
    \item Выделить входные параметры.
    \item Для каждого параметра определить допустимые и недопустимые области.
    \item Разбить каждую область на классы эквивалентности.
    \item Выбрать по одному или нескольку представителю из каждого класса.
    \item Сформировать тесты и ожидаемые результаты.
\end{enumerate}

\begin{example}
Функция принимает возраст пользователя $a$ и разрешает регистрацию, если:
\[
18 \leq a \leq 100.
\]

Классы эквивалентности:
\[
a < 18,\quad 18 \leq a \leq 100,\quad a > 100.
\]

Минимальные тесты:
\[
a = 17,\quad a = 25,\quad a = 101.
\]
\end{example}

\subsection{Анализ граничных значений}

Ошибки часто возникают на границах диапазонов, поэтому тестируются значения около границ.

\paragraph{Алгоритм}

\begin{enumerate}
    \item Найти границы допустимых диапазонов.
    \item Для каждой границы $b$ проверить значения $b-1$, $b$, $b+1$.
    \item Для вещественных величин использовать ближайшие представимые значения или малое $\varepsilon$.
\end{enumerate}

\begin{example}
Для диапазона:
\[
18 \leq a \leq 100
\]
проверяются:
\[
17,\ 18,\ 19,\ 99,\ 100,\ 101.
\]
\end{example}

\subsection{Таблицы решений}

Таблицы решений применяются, когда поведение системы зависит от комбинации условий.

\begin{example}
Пусть скидка предоставляется, если клиент является постоянным или сумма заказа не меньше $10000$.

\begin{center}
\begin{tabular}{c c c}
\toprule
Постоянный клиент & Сумма $\geq 10000$ & Скидка \\
\midrule
0 & 0 & 0 \\
0 & 1 & 1 \\
1 & 0 & 1 \\
1 & 1 & 1 \\
\bottomrule
\end{tabular}
\end{center}
\end{example}

\paragraph{Алгоритм}

\begin{enumerate}
    \item Выделить условия.
    \item Выделить действия.
    \item Построить все комбинации условий.
    \item Исключить невозможные комбинации.
    \item Для каждой оставшейся комбинации определить ожидаемое действие.
    \item Создать тесты.
\end{enumerate}

\subsection{Тестирование переходов состояний}

Метод применяется для систем, поведение которых зависит от состояния: протоколы, банкоматы, заказы, авторизация, жизненный цикл задач.

\begin{definition}
\textbf{Конечный автомат} --- модель $(S, E, \delta, s_0)$, где $S$ --- множество состояний, $E$ --- множество событий, $\delta$ --- функция переходов, $s_0$ --- начальное состояние.
\end{definition}

\begin{example}
Система авторизации:

\[
\text{Не авторизован}
\xrightarrow{\text{верный пароль}}
\text{Авторизован}
\xrightarrow{\text{выход}}
\text{Не авторизован}.
\]

При трёх неверных паролях:

\[
\text{Не авторизован}
\xrightarrow{\text{3 ошибки}}
\text{Заблокирован}.
\]
\end{example}

\paragraph{Алгоритм}

\begin{enumerate}
    \item Определить состояния.
    \item Определить события.
    \item Построить таблицу или граф переходов.
    \item Выбрать критерий покрытия: все состояния, все переходы, все пары переходов.
    \item Сгенерировать тестовые последовательности.
\end{enumerate}

\subsection{Тестирование сценариев использования}

Основано на пользовательских сценариях.

Применяется для:
\begin{itemize}
    \item системного тестирования;
    \item приёмочного тестирования;
    \item тестирования пользовательского интерфейса;
    \item проверки бизнес-процессов.
\end{itemize}

\begin{example}
Сценарий покупки товара:
\begin{enumerate}
    \item пользователь открывает каталог;
    \item выбирает товар;
    \item добавляет товар в корзину;
    \item оформляет заказ;
    \item оплачивает заказ;
    \item получает подтверждение.
\end{enumerate}
\end{example}

\section{Структурное тестирование и покрытие кода}

\subsection{Граф потока управления}

\begin{definition}
\textbf{Граф потока управления} --- ориентированный граф, вершины которого соответствуют базовым блокам программы, а рёбра --- возможным переходам управления.
\end{definition}

\begin{definition}
\textbf{Базовый блок} --- последовательность инструкций, в которую управление входит только в начало и выходит только из конца.
\end{definition}

\subsection{Критерии покрытия}

\begin{center}
\begin{tabular}{p{4cm} p{9cm}}
\toprule
\textbf{Критерий} & \textbf{Смысл} \\
\midrule
Покрытие операторов & Каждый оператор выполнен хотя бы один раз. \\
Покрытие ветвей & Каждая ветвь каждого условия выполнена хотя бы один раз. \\
Покрытие условий & Каждое простое логическое условие принимало значения true и false. \\
Покрытие путей & Каждый путь в графе управления выполнен хотя бы один раз. \\
MC/DC & Каждое атомарное условие независимо влияло на результат решения. \\
\bottomrule
\end{tabular}
\end{center}

\subsection{Пример покрытия}

Рассмотрим функцию:

\begin{lstlisting}[language=Python]
def f(x, y):
    if x > 0 and y > 0:
        return 1
    else:
        return 0
\end{lstlisting}

Для покрытия операторов достаточно теста:
\[
(x,y)=(1,1).
\]

Для покрытия ветвей нужны:
\[
(1,1),\quad (0,1).
\]

Для покрытия условий желательно проверить, что каждое из условий $x>0$ и $y>0$ принимает оба значения:
\[
(1,1),\quad (0,1),\quad (1,0).
\]

\section{Цикломатическая сложность}

\subsection{Определение}

\begin{definition}
\textbf{Цикломатическая сложность} программы --- метрика, оценивающая число линейно независимых путей в графе потока управления.
\end{definition}

Для связного графа потока управления:
\[
V(G) = E - N + 2,
\]
где:
\begin{itemize}
    \item $E$ --- число рёбер;
    \item $N$ --- число вершин;
    \item $V(G)$ --- цикломатическая сложность.
\end{itemize}

В более общем случае, если граф имеет $P$ компонент связности:
\[
V(G) = E - N + 2P.
\]

Для одной функции часто используется приближённая формула:
\[
V(G) = D + 1,
\]
где $D$ --- число точек принятия решений: \texttt{if}, \texttt{while}, \texttt{for}, \texttt{case}, логические разветвления.

\subsection{Теорема о цикломатической сложности}

\begin{theorem}
Пусть $G$ --- связный граф потока управления программы с одним входом и одним выходом. Если $N$ --- число вершин графа, а $E$ --- число рёбер, то число линейно независимых циклов в соответствующем графе равно:
\[
V(G)=E-N+2.
\]
Эта величина равна числу линейно независимых путей, которые необходимо протестировать для базисного тестирования путей.
\end{theorem}

\begin{proof}
Рассмотрим связный неориентированный граф, полученный из графа потока управления забыванием ориентации рёбер. Из теории графов известно, что размерность пространства циклов связного графа равна:
\[
\mu = E - N + 1.
\]

Докажем это.

Пусть $T$ --- остовное дерево графа $G$. Так как $T$ содержит все $N$ вершин и является деревом, число его рёбер равно:
\[
E_T = N - 1.
\]

Все остальные рёбра графа, не входящие в остовное дерево, называются хордами. Их число:
\[
E - E_T = E - (N-1)=E-N+1.
\]

Добавление каждой хорды к остовному дереву образует ровно один простой цикл, потому что в дереве между любыми двумя вершинами существует единственный простой путь. Если хорда соединяет вершины $u$ и $v$, то вместе с единственным путём из $u$ в $v$ по дереву она образует цикл.

Каждая такая хорда порождает один независимый цикл, поскольку соответствующий цикл содержит хорду, не входящую в циклы, построенные по другим хордам. Следовательно, число независимых циклов равно:
\[
\mu = E-N+1.
\]

В графе потока управления программы с одним входом и одним выходом для анализа путей обычно добавляют фиктивное ребро из выхода во вход. Это замыкает граф и превращает независимые пути программы в циклы замкнутого графа.

После добавления фиктивного ребра число рёбер становится $E+1$, а число вершин остаётся $N$. Поэтому размерность пространства циклов равна:
\[
(E+1)-N+1 = E-N+2.
\]

Следовательно:
\[
V(G)=E-N+2.
\]

Эта величина задаёт число базисных независимых путей: любой другой путь можно представить как комбинацию базисных путей с точки зрения прохождения рёбер графа. Теорема доказана.
\end{proof}

\subsection{Пример вычисления цикломатической сложности}

\begin{lstlisting}[language=Python]
def abs_max(a, b):
    if abs(a) > abs(b):
        return abs(a)
    else:
        return abs(b)
\end{lstlisting}

Есть одна точка ветвления:
\[
D=1.
\]

Следовательно:
\[
V(G)=D+1=2.
\]

Нужны два независимых теста:
\[
(a,b)=(5,3),\quad (a,b)=(2,-7).
\]

\section{Оценивание сложности программ}

\subsection{Асимптотическая сложность}

\begin{definition}
Пусть $f(n)$ и $g(n)$ --- неотрицательные функции. Говорят, что:
\[
f(n)=O(g(n)),
\]
если существуют положительные константы $C$ и $n_0$ такие, что для всех $n\geq n_0$:
\[
f(n)\leq Cg(n).
\]
\end{definition}

Также используются обозначения:

\[
f(n)=\Omega(g(n))
\]
--- нижняя асимптотическая оценка;

\[
f(n)=\Theta(g(n))
\]
--- точная асимптотическая оценка;

\[
f(n)=o(g(n))
\]
--- строго меньший порядок роста.

\subsection{Временная и пространственная сложность}

\begin{definition}
\textbf{Временная сложность} --- функция, оценивающая количество элементарных операций алгоритма в зависимости от размера входа.
\end{definition}

\begin{definition}
\textbf{Пространственная сложность} --- функция, оценивающая объём дополнительной памяти, необходимой алгоритму.
\end{definition}

\subsection{Пример анализа сложности}

\begin{lstlisting}[language=Python]
def contains_duplicate(a):
    for i in range(len(a)):
        for j in range(i + 1, len(a)):
            if a[i] == a[j]:
                return True
    return False
\end{lstlisting}

В худшем случае все пары элементов будут проверены:

\[
\sum_{i=0}^{n-1}(n-i-1)
=
(n-1)+(n-2)+\dots+1+0
=
\frac{n(n-1)}{2}.
\]

Следовательно:
\[
T(n)=\Theta(n^2).
\]

Дополнительная память:
\[
S(n)=\Theta(1).
\]

Оптимизированная версия:

\begin{lstlisting}[language=Python]
def contains_duplicate(a):
    seen = set()
    for x in a:
        if x in seen:
            return True
        seen.add(x)
    return False
\end{lstlisting}

В среднем:
\[
T(n)=\Theta(n),
\qquad
S(n)=\Theta(n).
\]

\subsection{Связь сложности и тестируемости}

Сложность программы влияет на тестирование:

\begin{itemize}
    \item высокая цикломатическая сложность увеличивает число независимых путей;
    \item вложенные условия затрудняют анализ граничных случаев;
    \item большое число параметров приводит к комбинаторному взрыву тестов;
    \item высокая связность модулей усложняет модульное тестирование;
    \item недетерминизм затрудняет воспроизводимость дефектов.
\end{itemize}

\section{Верификация программ}

\subsection{Динамическая и статическая верификация}

\begin{definition}
\textbf{Динамическая верификация} --- проверка программы при её выполнении. Основной пример --- тестирование.
\end{definition}

\begin{definition}
\textbf{Статическая верификация} --- проверка программы без её выполнения.
\end{definition}

К статическим методам относятся:

\begin{itemize}
    \item инспекции кода;
    \item статический анализ;
    \item проверка типов;
    \item формальная верификация;
    \item model checking;
    \item доказательство корректности программ.
\end{itemize}

\subsection{Контракты и утверждения}

Контракт функции обычно содержит:

\begin{itemize}
    \item \textbf{предусловие} --- что должно быть истинно перед вызовом;
    \item \textbf{постусловие} --- что должно быть истинно после завершения;
    \item \textbf{инвариант} --- свойство, сохраняющееся в определённых точках программы.
\end{itemize}

\begin{example}
Для функции деления:

\[
z = \frac{x}{y}
\]

предусловие:
\[
y \neq 0.
\]

постусловие:
\[
z \cdot y = x
\]
для точной арифметики.
\end{example}

\subsection{Тройка Хоара}

\begin{definition}
\textbf{Тройка Хоара} имеет вид:
\[
\{P\}\ C\ \{Q\},
\]
где $P$ --- предусловие, $C$ --- программа или команда, $Q$ --- постусловие.
\end{definition}

Смысл:
если перед выполнением команды $C$ истинно $P$ и команда завершается, то после её завершения истинно $Q$.

\subsection{Правило присваивания}

\begin{theorem}
Для команды присваивания $x := E$ корректно правило:
\[
\{Q[x \leftarrow E]\}\ x := E\ \{Q\},
\]
где $Q[x \leftarrow E]$ означает формулу $Q$, в которой все свободные вхождения $x$ заменены выражением $E$.
\end{theorem}

\begin{proof}
Пусть до выполнения присваивания истинна формула $Q[x \leftarrow E]$.

Команда $x := E$ вычисляет выражение $E$ в старом состоянии и присваивает полученное значение переменной $x$. После присваивания новое значение $x$ равно старому значению выражения $E$.

Так как до присваивания была истинна формула, полученная из $Q$ заменой $x$ на $E$, то после присваивания, когда $x$ стал равен этому значению $E$, формула $Q$ становится истинной.

Следовательно:
\[
\{Q[x \leftarrow E]\}\ x := E\ \{Q\}.
\]
Теорема доказана.
\end{proof}

\begin{example}
Нужно доказать:
\[
\{x+1>0\}\ y:=x+1\ \{y>0\}.
\]

Здесь:
\[
Q \equiv y>0.
\]

После замены $y$ на $x+1$:
\[
Q[y \leftarrow x+1] \equiv x+1>0.
\]

Следовательно, тройка корректна.
\end{example}

\section{Генерация тестов}

\subsection{Задача генерации тестов}

\begin{definition}
\textbf{Генерация тестов} --- это автоматическое или полуавтоматическое построение входных данных и ожидаемых результатов для проверки программы.
\end{definition}

Генерация тестов может быть:

\begin{itemize}
    \item ручной;
    \item случайной;
    \item комбинаторной;
    \item основанной на моделях;
    \item основанной на спецификациях;
    \item символической;
    \item мутационной;
    \item property-based.
\end{itemize}

\subsection{Случайная генерация тестов}

\paragraph{Алгоритм}

\begin{enumerate}
    \item Определить область входных данных.
    \item Задать распределение вероятностей.
    \item Сгенерировать случайный вход.
    \item Выполнить программу.
    \item Проверить результат с помощью оракула.
    \item Повторить заданное число раз или до обнаружения ошибки.
\end{enumerate}

\begin{example}
Для функции сортировки можно генерировать случайные массивы чисел.

Оракул:
\begin{enumerate}
    \item результат должен быть отсортирован;
    \item результат должен содержать те же элементы, что и исходный массив.
\end{enumerate}
\end{example}

\begin{lstlisting}[language=Python]
def is_sorted(a):
    return all(a[i] <= a[i+1] for i in range(len(a)-1))

def same_multiset(a, b):
    return sorted(a) == sorted(b)
\end{lstlisting}

\subsection{Property-based testing}

\begin{definition}
\textbf{Property-based testing} --- подход, при котором тесты генерируются автоматически, а корректность проверяется через общие свойства результата.
\end{definition}

Примеры свойств:

\begin{itemize}
    \item сортировка не меняет множество элементов;
    \item повторная сортировка не меняет результат:
    \[
    sort(sort(a)) = sort(a);
    \]
    \item для разворота списка:
    \[
    reverse(reverse(a))=a.
    \]
\end{itemize}

\subsection{Комбинаторная генерация тестов}

Если система имеет параметры:
\[
p_1,p_2,\dots,p_n,
\]
и каждый параметр имеет несколько значений, полный перебор может быть слишком большим.

\begin{example}
Пусть есть $5$ параметров по $4$ значения. Полный перебор:
\[
4^5=1024
\]
тестов.
\end{example}

Комбинаторное тестирование строит меньший набор тестов, покрывающий все комбинации значений параметров некоторого размера $t$.

\begin{definition}
\textbf{$t$-wise тестирование} --- метод, при котором каждая комбинация значений любых $t$ параметров встречается хотя бы в одном тесте.
\end{definition}

Частные случаи:

\begin{itemize}
    \item $t=1$ --- покрытие всех отдельных значений;
    \item $t=2$ --- попарное тестирование;
    \item $t=3$ --- покрытие всех троек параметров.
\end{itemize}

\subsection{Попарное тестирование}

\paragraph{Идея}

Многие дефекты проявляются из-за взаимодействия небольшого числа параметров. Попарное тестирование покрывает все пары значений параметров.

\begin{example}
Пусть есть параметры:

\[
A=\{0,1\},\quad B=\{0,1\},\quad C=\{0,1\}.
\]

Полный перебор даёт $8$ тестов. Попарное покрытие можно получить четырьмя тестами:

\[
\begin{array}{c c c}
A & B & C\\
\hline
0 & 0 & 0\\
0 & 1 & 1\\
1 & 0 & 1\\
1 & 1 & 0
\end{array}
\]

Пары $(A,B)$, $(A,C)$ и $(B,C)$ покрыты полностью.
\end{example}

\subsection{Упрощённый алгоритм жадной комбинаторной генерации}

\begin{enumerate}
    \item Построить множество всех требуемых $t$-комбинаций.
    \item Пока есть непокрытые комбинации:
    \begin{enumerate}
        \item сгенерировать множество кандидатов-тестов;
        \item для каждого кандидата посчитать, сколько непокрытых комбинаций он покрывает;
        \item выбрать лучший кандидат;
        \item добавить его в набор тестов;
        \item удалить покрытые комбинации.
    \end{enumerate}
\end{enumerate}

\begin{example}
Для параметров $A,B,C$ по два значения все пары:
\[
(A=0,B=0),\ (A=0,B=1),\dots
\]

Кандидат $(0,0,0)$ покрывает:
\[
(A=0,B=0),\quad (A=0,C=0),\quad (B=0,C=0).
\]

Затем выбирается следующий кандидат, покрывающий максимум ещё не покрытых пар.
\end{example}

\subsection{Генерация тестов по моделям}

\begin{definition}
\textbf{Model-based testing} --- генерация тестов на основе модели поведения системы: автомата, диаграммы состояний, спецификации протокола.
\end{definition}

\paragraph{Алгоритм}

\begin{enumerate}
    \item Построить модель системы.
    \item Выбрать критерий покрытия модели.
    \item Сгенерировать последовательности действий.
    \item Преобразовать абстрактные действия в реальные тестовые шаги.
    \item Выполнить тесты.
    \item Сравнить фактическое поведение с ожидаемым.
\end{enumerate}

\subsection{Символическое выполнение}

\begin{definition}
\textbf{Символическое выполнение} --- метод анализа программы, при котором вместо конкретных входных значений используются символические переменные.
\end{definition}

Для каждого пути строится условие пути.

\begin{example}
Рассмотрим программу:

\begin{lstlisting}[language=Python]
def f(x):
    if x > 10:
        return 1
    else:
        return 0
\end{lstlisting}

Пути:

\[
x>10 \Rightarrow return\ 1,
\]
\[
x\leq 10 \Rightarrow return\ 0.
\]

Тесты можно получить решением ограничений:
\[
x=11,\quad x=10.
\]
\end{example}

\paragraph{Алгоритм символической генерации тестов}

\begin{enumerate}
    \item Заменить входы символическими переменными.
    \item Интерпретировать программу над символическими выражениями.
    \item При каждом ветвлении добавлять условие в формулу пути.
    \item Для каждого достижимого пути решить систему ограничений.
    \item Получить конкретные входы.
    \item Выполнить программу на этих входах.
\end{enumerate}

\subsection{Concolic testing}

\begin{definition}
\textbf{Concolic testing} --- комбинированное конкретно-символическое выполнение, при котором программа выполняется на конкретном входе, но параллельно собираются символические ограничения пути.
\end{definition}

\paragraph{Алгоритм}

\begin{enumerate}
    \item Выбрать начальный конкретный вход.
    \item Выполнить программу.
    \item Собрать условия пройденного пути.
    \item Инвертировать одно из условий пути.
    \item Решить новую систему ограничений.
    \item Получить новый тест.
    \item Повторять до достижения нужного покрытия или исчерпания ресурсов.
\end{enumerate}

\begin{example}
Пусть путь задаётся условиями:
\[
x>0,\quad y=x+1,\quad y<5.
\]

Чтобы получить новый путь, можно заменить последнее условие:
\[
x>0,\quad y=x+1,\quad y\geq 5.
\]

Решение:
\[
x=4,\quad y=5.
\]
\end{example}

\subsection{Мутационное тестирование}

\begin{definition}
\textbf{Мутационное тестирование} --- метод оценки качества тестов, при котором в программу вносятся небольшие искусственные изменения, называемые мутантами.
\end{definition}

Примеры мутаций:

\begin{itemize}
    \item замена \texttt{>} на \texttt{>=};
    \item замена \texttt{+} на \texttt{-};
    \item удаление оператора;
    \item замена константы.
\end{itemize}

Если тесты обнаруживают отличие мутанта от исходной программы, говорят, что мутант ``убит''.

\[
Mutation\ Score =
\frac{\text{число убитых мутантов}}
{\text{число всех неэквивалентных мутантов}}.
\]

\begin{example}
Исходная функция:

\begin{lstlisting}[language=Python]
def max2(a, b):
    if a > b:
        return a
    return b
\end{lstlisting}

Мутант:

\begin{lstlisting}[language=Python]
def max2(a, b):
    if a >= b:
        return a
    return b
\end{lstlisting}

Для чисел результат часто совпадает. Такой мутант может быть эквивалентным относительно спецификации максимума двух чисел, потому что при $a=b$ оба варианта возвращают одинаковое значение.
\end{example}

\section{Системы генерации тестов}

\subsection{Классы инструментов}

\begin{center}
\begin{tabular}{p{4cm} p{9cm}}
\toprule
\textbf{Класс инструментов} & \textbf{Назначение} \\
\midrule
Unit testing frameworks & Запуск модульных тестов. \\
Property-based testing tools & Генерация случайных данных по свойствам. \\
Combinatorial testing tools & Генерация $t$-wise наборов. \\
Fuzzers & Генерация случайных или направленных входов для поиска сбоев. \\
Symbolic execution tools & Генерация тестов через решение ограничений. \\
Model-based testing tools & Генерация тестов по моделям. \\
Mutation testing tools & Оценка силы тестового набора. \\
\bottomrule
\end{tabular}
\end{center}

\subsection{Примеры систем}

\begin{itemize}
    \item \textbf{JUnit, pytest, GoogleTest} --- фреймворки модульного тестирования.
    \item \textbf{Hypothesis, QuickCheck} --- property-based testing.
    \item \textbf{NIST ACTS} --- комбинаторная генерация $t$-wise тестов.
    \item \textbf{AFL, libFuzzer, Honggfuzz} --- фаззинг.
    \item \textbf{KLEE} --- символическое выполнение программ на C.
    \item \textbf{Stryker, PIT} --- мутационное тестирование.
    \item \textbf{Tcases} --- генерация тестов по входным моделям.
\end{itemize}

\subsection{Фаззинг}

\begin{definition}
\textbf{Фаззинг} --- автоматическое тестирование программы большим количеством случайных, мутированных или специально сгенерированных входов.
\end{definition}

Фаззинг особенно полезен для:

\begin{itemize}
    \item парсеров;
    \item компиляторов;
    \item сетевых протоколов;
    \item библиотек обработки изображений, архивов, документов;
    \item поиска уязвимостей безопасности.
\end{itemize}

\paragraph{Упрощённый алгоритм фаззинга}

\begin{enumerate}
    \item Подготовить начальный корпус входов.
    \item Выбрать вход из корпуса.
    \item Выполнить мутацию: вставка, удаление, замена байтов.
    \item Запустить программу на новом входе.
    \item Если найден новый путь выполнения, сохранить вход.
    \item Если найден сбой, сохранить минимальный воспроизводимый пример.
    \item Повторять до остановки.
\end{enumerate}

\begin{example}
Исходный вход:
\[
\texttt{12+34}
\]

Мутации:
\[
\texttt{12++34},\quad
\texttt{12/0},\quad
\texttt{((((12}
\]

Такие входы могут обнаружить ошибки парсера арифметических выражений.
\end{example}

\section{Нефункциональное тестирование}

\subsection{Тестирование производительности}

Проверяет время ответа, пропускную способность, масштабируемость и потребление ресурсов.

Виды:

\begin{itemize}
    \item \textbf{нагрузочное тестирование} --- поведение при ожидаемой нагрузке;
    \item \textbf{стресс-тестирование} --- поведение за пределами нормальной нагрузки;
    \item \textbf{тестирование масштабируемости} --- изменение производительности при увеличении ресурсов;
    \item \textbf{тестирование стабильности} --- длительная работа под нагрузкой.
\end{itemize}

\subsection{Тестирование безопасности}

Проверяет устойчивость к атакам:

\begin{itemize}
    \item SQL-инъекции;
    \item XSS;
    \item CSRF;
    \item обход авторизации;
    \item утечки данных;
    \item небезопасная десериализация;
    \item ошибки криптографического использования.
\end{itemize}

\subsection{Тестирование удобства использования}

Оценивает, насколько система понятна пользователю.

Проверяются:

\begin{itemize}
    \item простота интерфейса;
    \item понятность сообщений об ошибках;
    \item доступность;
    \item число действий для выполнения задачи;
    \item соответствие ожиданиям пользователя.
\end{itemize}

\subsection{Тестирование совместимости}

Проверяет работу системы в разных окружениях:

\begin{itemize}
    \item операционные системы;
    \item браузеры;
    \item версии библиотек;
    \item устройства;
    \item сетевые условия.
\end{itemize}

\section{Области применения методов тестирования}

\begin{center}
\begin{tabular}{p{4cm} p{8cm} p{3cm}}
\toprule
\textbf{Метод} & \textbf{Где применять} & \textbf{Типичные дефекты} \\
\midrule
Классы эквивалентности & Формы, API, валидация данных & Неверная обработка диапазонов \\
Граничные значения & Числовые диапазоны, длины строк, лимиты & Ошибки на границах \\
Таблицы решений & Бизнес-правила, тарифы, права доступа & Неверные комбинации условий \\
Переходы состояний & Протоколы, заказы, авторизация & Неверные состояния \\
Сценарии использования & Приёмочное и системное тестирование & Нарушение бизнес-процесса \\
Покрытие операторов & Модульное тестирование & Непроверенный код \\
Покрытие ветвей & Условная логика & Непроверенные ветви \\
MC/DC & Критические системы & Ошибки сложных условий \\
Фаззинг & Парсеры, протоколы, безопасность & Падения, переполнения, уязвимости \\
Комбинаторное тестирование & Конфигурации, совместимость & Ошибки взаимодействия параметров \\
Мутационное тестирование & Оценка качества тестов & Слабые проверки \\
\bottomrule
\end{tabular}
\end{center}

\section{Тестовый оракул}

\begin{definition}
\textbf{Тестовый оракул} --- механизм определения того, является ли результат выполнения теста правильным.
\end{definition}

Виды оракулов:

\begin{itemize}
    \item точный ожидаемый результат;
    \item сравнение с эталонной реализацией;
    \item проверка свойств;
    \item метаморфное тестирование;
    \item ручная проверка;
    \item проверка инвариантов.
\end{itemize}

\subsection{Метаморфное тестирование}

\begin{definition}
\textbf{Метаморфное тестирование} --- подход, при котором проверяется не конкретный результат, а отношение между результатами нескольких запусков.
\end{definition}

\begin{example}
Для функции сортировки:
\[
sort(a) = sort(reverse(a)).
\]

Если это свойство нарушено, значит в сортировке есть ошибка.
\end{example}

\section{Минимальный пример комплексного тестирования}

Пусть требуется протестировать функцию:

\begin{lstlisting}[language=Python]
def category(age):
    if age < 0:
        return "invalid"
    if age < 18:
        return "child"
    if age <= 65:
        return "adult"
    return "senior"
\end{lstlisting}

\subsection{Классы эквивалентности}

\[
age < 0,\quad 0 \leq age < 18,\quad 18 \leq age \leq 65,\quad age > 65.
\]

Тесты:
\[
-1,\quad 10,\quad 30,\quad 70.
\]

\subsection{Граничные значения}

Границы:
\[
0,\quad 18,\quad 65.
\]

Тесты:
\[
-1,\ 0,\ 1,\ 17,\ 18,\ 19,\ 64,\ 65,\ 66.
\]

\subsection{Цикломатическая сложность}

Три условия:
\[
D=3.
\]

Следовательно:
\[
V(G)=D+1=4.
\]

Нужно минимум четыре независимых пути:

\[
age<0,
\]
\[
0\leq age<18,
\]
\[
18\leq age\leq65,
\]
\[
age>65.
\]

\section{Заключение}

Отладка, тестирование, верификация и оценивание сложности образуют единый комплекс методов обеспечения качества программного обеспечения.

\begin{itemize}
    \item Отладка направлена на нахождение и устранение причин уже обнаруженных дефектов.
    \item Тестирование позволяет находить дефекты и оценивать качество программы.
    \item Верификация проверяет соответствие программы спецификации.
    \item Оценивание сложности помогает прогнозировать трудность сопровождения и тестирования.
    \item Генерация тестов позволяет автоматизировать построение тестовых данных и бороться с комбинаторным ростом вариантов.
\end{itemize}

На практике методы комбинируются: функциональные тесты проверяют требования, структурные тесты контролируют покрытие кода, статический анализ выявляет дефекты без запуска, фаззинг ищет нестандартные сбои, а формальная верификация применяется для наиболее критичных компонентов.

\end{document}
